\documentclass[12pt]{article}
\usepackage[utf8]{inputenc}
\usepackage[T1]{fontenc}
\usepackage{hyperref}
\title{Molo Language Resources for K-12 Curriculum Development}
\author{Compiled by Claude Sonnet 5.5 (claude-sonnet-5-5) for Andr\'as Kornai}
\date{}
\begin{document}
\maketitle

\section{Introduction}
Molo (ISO 639-3 zmo, Glottocode molo1257) is a Nilo-Saharan language of the Eastern Jebel group in Sudan. It is spoken near Jebel Malkan, close to the Ethiopian border. It is also called Malkan or Tura-Ka-Molo. UNESCO rates it as critically endangered, with about 100 speakers according to 1988 data. Some sources suggest it may now be extinct. Search found only reference entries. It found no teaching materials.

\section{Primary Educational Resources}
None were found. The language is not known to be taught in schools.

\section{Linguistic and Reference Materials}
\begin{itemize}
\item Wikipedia, Molo language: \url{https://en.wikipedia.org/wiki/Molo_language}
\item Wikipedia, Eastern Jebel languages: \url{https://en.wikipedia.org/wiki/Eastern_Jebel_languages}
\item ASJP, Molo word list: \url{https://asjp.clld.org/languages/MOLO}
\end{itemize}

\section{Religious and Community Resources}
None were found.

\section{Digital and Audio Resources}
\begin{itemize}
\item Endangered Languages Project, Molo entry: \url{https://endangeredlanguages.com/lang/3134}
\end{itemize}
I did not check whether this entry links any samples.

\section{Curriculum Development Context}
With so few and such elderly speakers, the realistic goal is documentation and revitalisation rather than a full K-12 curriculum. The related Gaam language has many more speakers. Any Molo programme would need first to confirm whether speakers remain.

\end{document}